\section{Methods}\label{methods}

\subsection{Data Sources and Preparation}\label{data-sources-and-preparation}

To assess models, we assembled two datasets suitable for forecasting opioid-related fatal overdoses annually at the census tract level. Our relatively coarse annual temporal scale was chosen to match the frequency at which decision-makers might set new priorities and at which new reliable data become available. We chose the census tract spatial scale due to its potential for targeted interventions at a sub-municipality level. Each census tract by design contains a mean count of 4000 people 
(with a range of 1200-8000 \citep{uscensusbureauGlossary2022}
) . For many (but not all) interventions, costs scale with population size, and thus the cost of deploying an intervention in any tract is roughly uniform.

\subsubsection{\texorpdfstring{Data source 1: Cook County, Illinois }{Data source 1: Cook County, Illinois }}\label{data-source-1-cook-county-illinois}

We obtained fully de-identified data from the Cook County Government Medical Examiner Case Archive\citep{MedicalExaminerCase} for opioid-involved overdose deaths from August 2014 (the first date records are available) to May 2023. These data contained every fatal incident under the medical examiner's jurisdiction that was determined to have any opioid as a primary cause. We used the provided incident latitude and longitude to map each overdose fatality to one of 1328 census tracts. Because the underlying data is in the public domain, we make our processed Cook County data available in our shared repository.

\subsubsection{Data source 2: Massachusetts}\label{data-source-2-massachusetts}

We obtained death certificate data from the Massachusetts Registry of Vital Records and Statistics for opioid-involved overdose deaths between 2001 and 2019. These deaths were defined as unintentional, intentional, and undetermined drug poisonings containing an opioid code (ICD-10 codes T.40.0-T40.4, or T40.6) as a ``multiple cause of death''. Each fatal overdose is linked to a calendar date as well as a residential street address. Decedent addresses for the place of residence at the time of the fatal overdose were geocoded, assigning latitude and longitude measures to each event that was then mapped to one of 1620 census tracts using the 2020 census tract boundaries.

\subsubsection{Dataset Preparation}\label{dataset-preparation}

For each dataset, we computed the observed number of fatal overdose events \(y_{s,t}\) at time unit \(t\) for individuals residing in spatial tract \(s\). 
We employed open tools \citep{freemanCensusgeocodeThinPython}
\ that utilize the US Census Geocoding API to map locations (street address or latitude/longitude) to its corresponding census tract, using the tract boundaries for the states of Massachusetts and Illinois defined by the U.S. Census Bureau in 2020.

In each dataset, a uniform set of covariates is available as input for prediction models. At each time period \emph{t} and spatial region \emph{s}, we provide the history of fatal overdose counts from previous times in that region, as well as the spatial location (numerical latitude and longitude of the tract), and timestamp (numerical time, measured in years since the first available year for that dataset). Further, for each census tract \emph{s} at time \emph{t}, an optional covariate vector represents social vulnerability across socioeconomic status, age-related demographics, minority status, housing, and composite dimensions using percentile ranking across all census tracts in that state. 
These features stem from the five dimensions of the Social Vulnerability Index (SVI)\citep{CDCATSDRSVI2022}
\ published between 2000 and 2018 for every U.S. state. Published tract values are updated every five years; we selected the closest value to each time period t. These SVI features were chosen for their simplicity and portability, mirroring the role of higher-dimensional socioeconomic covariates in previous studies \citep{bauerSmallAreaForecasting2023,marksIdentifyingCountiesRisk2021}.

\subsection{Metrics}\label{metrics}

To evaluate model forecasting accuracy against observed mortality in a test period, suitable performance metrics were essential. Our study considered both commonly-used error-based metrics and a new intervention-focused metric.

\subsubsection{Error-based metrics}\label{error-based-metrics}

Model performance is often assessed via summary statistics of the errors between predicted and observed mortality across all S spatial regions in the test period. Within this category, two common metrics are Root Mean Squared Error (RMSE) and Mean Absolute Error (MAE), both defined in Equation 1 below. RMSE calculates the square root of the average squared errors, while MAE computes the average of absolute errors.

% FIXME: RMSE formula did not port from Pandoc -- only the MAE half is present below; restore the RMSE expression (sqrt of mean squared errors) before the comma
\(RMSE\  = \ ,\ \ \ \ \ MAE\  = \frac{1}{S}\sum_{s = 1}^{S}|y_{s} - \widehat{y_{s}}|\) (1)

Both RMSE and MAE have been concretely used as the primary metrics to assess opioid overdose forecasting\citep{ertugrulCASTNetCommunityAttentiveSpatioTemporal2019,bauerSmallAreaForecasting2023, sumetskySpatiotemporalModelingOpioid2017}
. RMSE can be more sensitive to large errors due to its use of squaring.

\subsubsection{Intervention-focused metric}\label{intervention-focused-metric}

In our intended use case of mitigating the opioid crisis, stakeholders at a public health agency could use a forecasting model to select a targeted subset of all possible census tracts in which to deploy an intervention in the near future. We assume these actors have a limited budget, allowing intervention deployment in a maximum of K of the S regions in their jurisdiction. For a given model, we can obtain its recommended set of K regions, which we refer to as the intervention set \emph{I}, in two steps. Step 1: predict mortality counts for all S regions in the test period. Step 2: identify the K regions with the K highest predictions (breaking ties at random), and store these as the recommended set \emph{I.}

To evaluate such intervention recommendations, we need new metrics. The error-based metrics defined earlier (RMSE or MAE) aggregate error for all S regions equally. They do not directly measure if a forecast's guess of the K highest-risk areas specifically aligns with the actual areas with highest mortality. We thus wish to design a metric better aligned with how stakeholders will determine and assess intervention priorities. We suggest a model is favorable if, during the test period, the total count of fatal overdose events in its recommended set \emph{I} of K tracts is as large as possible. This would indicate the model is good at identifying where adverse events will occur, and thus increase the possibility that stakeholders could ``reach'' and hopefully mitigate these events via interventions targeted at the model's recommended tracts.

Our proposed metric, ``best possible reach'' (BPR), assesses a model's recommendations via a ratio of two numbers. First, the numerator counts how many opioid-related fatalities actually occurred in the model's recommended set \emph{I} of size K. Second, the denominator counts the total number of opioid-related fatalities in the K regions that would be chosen with perfect hindsight of the actual count vector y = {[}y\textsubscript{1}, \ldots, y\textsubscript{S}{]} of fatalities in the test period. Mathematically, we define BPR as follows

\(BPR\  = \ \frac{\sum_{k \in I}^{}y_{k}}{\sum_{k \in TopKInds(y)}^{}y_{k}}\) = \(\frac{\ \#\ events\ in\ K\ regions\ picked\ by\ model}{\#\ events\ in\ actual\ K\ highest\ count\ regions}\)

where TopKInds(y) denotes a function that returns the distinct indices of the K largest elements of vector \emph{y}.

For public health applications, BPR holds a practical interpretation as the proportion of fatal overdose events the current model's interventions would reach, compared to the perfect foresight of future events. BPR's numerical value by definition will have a minimum of 0.0 and a maximum of 1.0. We typically convert the fractional BPR to a percentage ranging from 0-100\%, denoted as \%BPR. A higher \%BPR value signifies a more effective model at deciding where to intervene. A value of 100\% indicates perfect decision-making given a limited budget: there is no other set of K regions any model could have recommended that would reach more events.

Although independently developed by our team (see our preliminary workshop paper \citet{heutonZIGP2022}
), our proposed BPR metric closely resembles the metric suggested by a recent pre-registered trial\citep{marshallPreventingOverdoseUsing2022} and a later feasibility study\citep{allenTranslatingPredictiveAnalytics2023} to evaluate opioid overdose forecasts in Rhode Island. The primary distinction lies in the denominator: our BPR sums only the top K indices, while the alternative includes all S regions. We prefer our definition due to \%BPR's consistent range of 0-100\%, with 100\% representing a model that could not have made a better decision given the limited budget of K regions. In contrast, the alternative definition's maximum value fluctuates based on observed data in the test period, making it difficult to know if another model could have done better.

\subsection{Models}\label{models}

Below we define a range of possible forecasting methods that can be used for our where-to-intervene prediction task. All methods are trained and evaluated via a common protocol using the same provided splits (train/validation/test) of the two available datasets (Cook County and Massachusetts). Thorough details about model fitting and hyperparameter tuning are provided in the Supplement. This study was reviewed by the Tufts University Health Sciences Institutional Review Board and deemed to be non-human subjects research.

\subsubsection{Simple Baseline Models}\label{simple-baseline-models}

We study several easy-to-implement baseline models to highlight their comparative strengths. Public health practitioners seeking data-driven allocation of scarce intervention resources without sophisticated modeling could easily use these approaches.

Our first baseline, dubbed \emph{all zeroes,} predicts uniformly across all S tracts that zero fatal overdoses will occur in the test period. This model, by definition, ranks all tracts as equally high-risk, so for metrics like BPR that require a set of K recommended regions we report an average over many samples of K distinct regions selected uniformly at random.

Our second baseline, known as \emph{last year,} predicts that the mortality at the next year in a specific location will mirror the mortality observed at the most recent recorded year for that location.

The final baseline we consider is a \emph{historical average,} which predicts the next timestep's mortality as an average of all mortality counts observed over the preceding W timesteps. Here, the appropriate ``look back'' period length W serves as a model hyperparameter.

\subsubsection{Complex Models}\label{complex-models}

Next, we consider several more flexible models with parameters that can be fit to the data. The first is the \emph{weighted historical} average: a weighted average of the previous W years of overdose events. This is more flexible than \emph{historical average}, because each year's count is multiplied by a customized weight coefficient.

The next model we consider is a Generalized Linear Model (GLM) with a Poisson likelihood. This model assumes that fatal overdose count \emph{y} for spatial tract \emph{s} at time \emph{t} is modeled by a Poisson distribution where the log of the mean parameter is a linear function of the covariate vector \emph{x} for that tract and time:

\(y|\ x,\bm{\theta}\  \sim \ Poisson(\lambda\  = e^{\bm{\theta}^{T}x})\ \)

GLMs have limited flexibility due to the assumption of a monotonic relationship between covariates and fatal event count. We thus also include a Gradient Boosted Trees\citep{friedmanGreedyFunctionApproximation2001} model, a popular more flexible ensemble of regression trees. Previous studies of opioid forecasting \citep{allenTranslatingPredictiveAnalytics2023,schellIdentifyingPredictorsOpioid2022} have used similar tree ensembles.

In addition, we include three spatially-sophisticated statistical models used in recent opioid overdose forecasting applications. First, we include a \emph{Gaussian Process} model\citep{allenTranslatingPredictiveAnalytics2023} for its ability to flexibly capture spatial and temporal correlations. We use similar covariance functions (``kernels'') to prior overdose forecasting work (details in the supplement). Next, \emph{Bayesian Spatio-Temporal (BST)} models\citep{bauerSmallAreaForecasting2023} use a Markov Random Field to model inherent spatial and temporal trends. Thirdly, \emph{NBSpLag} denotes a negative binomial regression model with spatially lagged features\citep{marksIdentifyingCountiesRisk2021}, where each tract is informed by its spatial neighbors. In a variable selection experiment\citep{marksIdentifyingCountiesRisk2021}, these spatially lagged covariates were found to be the most predictive features. Unlike previous evaluations of both \emph{BST}\citep{bauerSmallAreaForecasting2023} and \emph{NBSpLag} models, our study compares to the rich set of baselines described above.

Finally, we include \emph{CASTNet}\citep{ertugrulCASTNetCommunityAttentiveSpatioTemporal2019}, a neural network approach custom developed for opioid-overdose forecasting. Unlike previous methods, CASTNet employs multi-head attentional networks that allow predictions at a given location to be informed by learned ``communities'' of regions.

\subsection{Experimental Protocol}\label{experimental-protocol}

We applied each of the models described above separately to the Cook County, IL dataset and the MA dataset. In each case, we sought to use available historical counts of opioid-related fatal overdoses (together with other covariates described above) to predict future fatal overdose counts in each census tract. We further assessed how these predictions can be used to recommend where to intervene in the near future.

For training on a dataset, for all S regions we assemble covariate vector, fatality count pairs \((x_{s,t},\ y_{s,t})\) for each year in the training set (t = 2010-2018 for Massachusetts, 2015-2019 for Cook County). The historical covariates inside each \(x_{s,t}\) vector can summarize the recent history of \(W\) previous years (W=10 for Massachusetts, W=5 for Cook County). Hyperparameters are chosen to maximize performance as assessed by BPR on a validation set of data from the year prior to evaluation (2019 for Massachusetts, 2020 for Cook County. Finally, models are evaluated on predictions for the final two years (2020-2021 in Massachusetts, 2021-2022 in Cook County).

From each model, we obtained predictions for each of the S tracts in each test year. We then computed each evaluation metric (RMSE, MAE, BPR) as well as an interval that quantifies our uncertainty in its precise value. Inspired by resampling methods for uncertainty quantification\citep{SpatialSamplingResampling}, for each test year we obtained 50 different without-replacement samples of 1370 of the 1620 locations in MA (1078 of the 1328 in Cook County, IL), and recorded all metrics of interest for each sample. Our reported intervals quantify the min-max range of these 50 samples. We chose the number of tracts retained in each catchment area so that each sample on average retains 85\% of all fatal overdose events.
